\section{Discussion and Conclusion}\label{sec:conclusion}

This project is intentionally narrow. Rather than attempting to reproduce the entire s1 training recipe, it concentrates on the part of the method that is easiest to inspect and validate locally: budget forcing at decoding time. That focus is appropriate for a course project because it yields a useful test of reproducibility without turning the project into an infrastructure exercise.

The practical lesson is that simple inference-time interventions are worth studying precisely because they are falsifiable. If budget forcing improves accuracy in the local setup, the report can analyze how much improvement is obtained per unit of extra compute and whether the gain appears stable across triggers. If it fails or produces mixed results, that outcome is still valuable because it clarifies the dependency of the original claim on model family, prompt format, hardware, and benchmark details.

Future work should study three directions. First, trigger selection should move beyond fixed phrases toward adaptive or learned continuation prompts. Second, anti-loop safeguards should be incorporated so that extra thinking budget is spent on verification rather than repetition. Third, budget forcing should be tested in settings where reasoning must be combined with external knowledge or tools, including retrieval-augmented pipelines and non-mathematical tasks.
\label{sec:conclusion}

This repository captures an instructive slice of the broader reasoning-model landscape. Its main value is not that it recreates every part of the original s1 pipeline, but that it operationalizes one of the paper's clearest ideas: test-time compute can be controlled directly at decoding time. That makes budget forcing a useful teaching and reproduction target because the core mechanism is understandable, inspectable, and experimentally modifiable.

At the same time, the project also highlights a recurring tension in contemporary LLM research. Methods advertised as simple may still depend on heavyweight runtime conditions that are invisible in the algorithm sketch. For that reason, a careful reproduction should treat environment notes, model substitutions, and execution blockers as first-class results rather than incidental details.

The final conclusion of this report should be updated once Agent A finishes the runtime validation and result generation. If the experiments succeed, the discussion should explain whether additional enforced thinking improved accuracy and how stable that improvement was across triggers. If the experiments are blocked, the conclusion should instead emphasize what the repository confirms at the implementation level and what remains unverified at the systems level.